\chapter{Traits and Impl Blocks}\label{ch:traits}
\index{trait}
\index{impl}

In \cref{ch:structs} we attached behaviour to structs through traits and
impl blocks, but we treated them as a recipe to follow without explaining
the ingredients. This chapter puts traits centre stage. We will see how to
declare interfaces, implement them for multiple types, compose impl blocks,
and use traits as the backbone of polymorphic design.

Think of a trait as a \emph{contract}---a promise that a type will provide
certain methods. An impl block is the \emph{fulfilment} of that promise. By
separating the ``what'' (the trait) from the ``how'' (the impl), Lattice
lets you write code that operates on \emph{any type that speaks a given
protocol}, without knowing the concrete type in advance.

%% ─────────────────────────────────────────────────────────────────────
\section{Declaring Traits}\label{sec:trait-declaring}
\index{trait!declaration}

A trait declaration lists method signatures with no bodies:

\begin{lstlisting}[language=Lattice, caption={A trait with a single method}]
trait Describable {
    fn describe(self: any) -> String
}
\end{lstlisting}

The keyword \lstinline|trait| is followed by a name (conventionally an
adjective or capability noun), then braces containing one or more method
signatures. Each signature specifies:

\begin{itemize}
  \item The method name.
  \item A parameter list. The first parameter is almost always
    \lstinline|self: any|, representing the value the method is called on.
  \item An optional return type.
\end{itemize}

\begin{defnbox}{Trait}
A \emph{trait} is a named set of method signatures that defines a
behavioural interface. It contains no implementation---only declarations
of the methods a type must provide.
\end{defnbox}

\subsection{Multi-Method Traits}

A trait can declare any number of methods:

\begin{lstlisting}[language=Lattice, caption={A trait with multiple methods}]
trait Collection {
    fn push_item(self: any, item: any) -> any
    fn pop_item(self: any) -> any
    fn peek_item(self: any) -> any
    fn is_empty(self: any) -> any
}
\end{lstlisting}

This \lstinline|Collection| trait declares four methods. Any type that
implements \lstinline|Collection| is expected to provide all four.

\subsection{Traits with Extra Parameters}

Methods are not limited to \lstinline|self|. They can accept additional
parameters:

\begin{lstlisting}[language=Lattice, caption={Traits with extra parameters}]
trait Scalable {
    fn scale(self: any, factor: any) -> any
}

trait Combinable {
    fn combine(self: any, other: any) -> any
}
\end{lstlisting}

The \lstinline|Scalable| trait expects a method that takes a scaling factor;
\lstinline|Combinable| expects a method that merges two values of the same
kind.

\begin{notebox}{Trait Bodies Are Signatures Only}
You cannot write method bodies inside a trait declaration. The trait is
purely a contract. All implementation goes in \lstinline|impl| blocks.
This is a deliberate design choice: it keeps the interface description
separate from the implementation, making both easier to read.
\end{notebox}

%% ─────────────────────────────────────────────────────────────────────
\section{Implementing Traits for Structs}\label{sec:trait-implementing}
\index{impl!block}
\index{trait!implementation}

An \lstinline|impl| block provides method bodies for a specific type:

\begin{lstlisting}[language=Lattice, caption={Implementing a trait}]
struct Animal { kind: any, name: any, legs: any }

trait Displayable {
    fn display(self: any) -> String
}

impl Displayable for Animal {
    fn display(self: any) -> String {
        return self.name + " the " + self.kind
            + " (" + to_string(self.legs) + " legs)"
    }
}

let dog = Animal { kind: "dog", name: "Rex", legs: 4 }
print(dog.display())  // Rex the dog (4 legs)
\end{lstlisting}

The syntax is \lstinline|impl TraitName for TypeName { ... }|, and inside
the braces you write full function definitions. Each method must match a
signature declared in the trait.

\subsection{How \texttt{self} Works}

The first parameter, \lstinline|self|, receives a \emph{copy} of the value
the method was called on. This follows Lattice's value semantics:

\begin{lstlisting}[language=Lattice, caption={Self is a copy}]
struct Counter { value: any }

trait Incrementable {
    fn increment(self: any) -> Counter
}

impl Incrementable for Counter {
    fn increment(self: any) -> Counter {
        return Counter { value: self.value + 1 }
    }
}

let c = Counter { value: 10 }
let c2 = c.increment()
print(c.value)   // 10 (unchanged)
print(c2.value)  // 11
\end{lstlisting}

Because \lstinline|self| is a copy, ``mutating'' methods must return a new
instance. The original remains untouched. This is the same pattern we saw
with struct methods in \cref{sec:struct-methods}.

\subsection{Accessing Fields through \texttt{self}}

Inside a method body, \lstinline|self| is a regular struct value. You read
its fields with dot notation:

\begin{lstlisting}[language=Lattice, caption={Reading fields from self}]
struct Rect { w: any, h: any }

trait Geometry {
    fn area(self: any) -> any
    fn perimeter(self: any) -> any
}

impl Geometry for Rect {
    fn area(self: any) -> any {
        return self.w * self.h
    }

    fn perimeter(self: any) -> any {
        return 2 * (self.w + self.h)
    }
}

let r = Rect { w: 5, h: 3 }
print(r.area())       // 15
print(r.perimeter())  // 16
\end{lstlisting}

\subsection{Method Chaining}

When a method returns a struct of the same type, you can chain calls:

\begin{lstlisting}[language=Lattice, caption={Method chaining}]
struct Stack { items: any, size: any }

trait Collection {
    fn push_item(self: any, item: any) -> any
    fn peek_item(self: any) -> any
}

impl Collection for Stack {
    fn push_item(self: any, item: any) -> any {
        flux new_items = []
        for i in 0..self.size {
            new_items.push(self.items[i])
        }
        new_items.push(item)
        return Stack { items: new_items, size: self.size + 1 }
    }

    fn peek_item(self: any) -> any {
        if self.size == 0 {
            return nil
        }
        return self.items[self.size - 1]
    }
}

let s = Stack { items: [], size: 0 }
let top = s.push_item(1).push_item(2).push_item(3).peek_item()
print(top)  // 3
\end{lstlisting}

Each \lstinline|push_item| returns a new \lstinline|Stack|, so we can
immediately call the next method on the result. The chain reads like a
pipeline: push 1, push 2, push 3, peek.

\begin{tipbox}{Chaining Reads Naturally}
Method chaining works because Lattice methods that return the same type
create a ``fluent interface.'' This is a powerful pattern for building up
data structures or transformation pipelines without intermediate variables.
\end{tipbox}

%% ─────────────────────────────────────────────────────────────────────
\section{Multiple Impl Blocks per Struct}\label{sec:trait-multiple-impl}
\index{impl!multiple blocks}

A struct can implement as many traits as it needs. Each trait gets its
own \lstinline|impl| block:

\begin{lstlisting}[language=Lattice, caption={Three traits on one struct}]
struct Point { x: any, y: any }

trait Describable {
    fn describe(self: any) -> String
}

trait HasArea {
    fn area(self: any) -> any
}

trait Scalable {
    fn scale(self: any, factor: any) -> any
}

impl Describable for Point {
    fn describe(self: any) -> String {
        return "Point(" + to_string(self.x) + ", " + to_string(self.y) + ")"
    }
}

impl HasArea for Point {
    fn area(self: any) -> any {
        return self.x * self.y
    }
}

impl Scalable for Point {
    fn scale(self: any, factor: any) -> any {
        return Point { x: self.x * factor, y: self.y * factor }
    }
}

let p = Point { x: 3, y: 4 }
print(p.describe())  // Point(3, 4)
print(p.area())      // 12

let q = p.scale(3)
print(q.x)  // 9
print(q.y)  // 12
\end{lstlisting}

The \lstinline|Point| struct now speaks three protocols:
\lstinline|Describable|, \lstinline|HasArea|, and \lstinline|Scalable|. Each
impl block is self-contained and can be placed anywhere in the same file
(or module).

\subsection{How Method Registration Works}

Under the hood, the compiler registers each method as a global function
with a composite key. When you write:

\begin{lstlisting}[language=Lattice]
impl HasArea for Rect {
    fn area(self: any) -> any { ... }
}
\end{lstlisting}

the compiler emits the method body as a function and stores it under the
key \texttt{Rect::area}. When you later call \lstinline|r.area()|, the VM:

\begin{enumerate}
  \item Checks for a callable field named \lstinline|area| on the struct
    (there is none).
  \item Looks up the global \texttt{Rect::area}.
  \item Calls it with a clone of \lstinline|r| as \lstinline|self|.
\end{enumerate}

You can see this in \filename{src/stackcompiler.c}: the \texttt{ITEM\_IMPL}
case builds the key string \lstinline|"TypeName::methodName"| and emits
\texttt{OP\_DEFINE\_GLOBAL} with that key. The VM's \texttt{OP\_INVOKE}
handler resolves the method at call time.

\begin{notebox}{Callable Fields Take Priority}
If a struct has both a closure field and a trait method with the same name,
the closure field wins. The runtime checks fields first, then falls through
to the impl registry. This is by design: it lets you override a trait
method on a per-instance basis if needed (see \cref{sec:struct-callable-fields}).
\end{notebox}

%% ─────────────────────────────────────────────────────────────────────
\section{Designing with Interfaces}\label{sec:trait-design}
\index{trait!design patterns}
\index{interface}

Traits are Lattice's interface mechanism. Effective use of traits leads to
code that is modular, testable, and extensible. Here are some patterns.

\subsection{The Same Trait on Different Types}

The real power of traits emerges when multiple types implement the same
interface. Functions that operate on the trait---rather than a specific
struct---work with \emph{any} implementing type:

\begin{lstlisting}[language=Lattice, caption={One trait, two types}]
struct Circle { radius: Float }
struct Rect { width: Float, height: Float }

trait Describable {
    fn describe(self: any) -> String
}

trait HasArea {
    fn area(self: any) -> Float
}

impl Describable for Circle {
    fn describe(self: any) -> String {
        return "Circle(r=" + to_string(self.radius) + ")"
    }
}

impl HasArea for Circle {
    fn area(self: any) -> Float {
        return 3.14159 * self.radius * self.radius
    }
}

impl Describable for Rect {
    fn describe(self: any) -> String {
        return "Rect(" + to_string(self.width) + "x" + to_string(self.height) + ")"
    }
}

impl HasArea for Rect {
    fn area(self: any) -> Float {
        return self.width * self.height
    }
}

let c = Circle { radius: 5.0 }
let r = Rect { width: 4.0, height: 6.0 }

print(c.describe())  // Circle(r=5)
print(c.area())      // 78.53975

print(r.describe())  // Rect(4x6)
print(r.area())      // 24
\end{lstlisting}

Both \lstinline|Circle| and \lstinline|Rect| implement \lstinline|Describable|
and \lstinline|HasArea|. The methods have different bodies---each type knows
how to describe and measure itself---but the \emph{interface} is identical.

\subsection{Traits for Separation of Concerns}

Suppose you are building a notification system. You might define separate
traits for each concern:

\begin{lstlisting}[language=Lattice, caption={Separating concerns with traits}]
struct EmailAlert {
    to: String,
    subject: String,
    body: String
}

trait Sendable {
    fn send(self: any) -> String
}

trait Loggable {
    fn log_entry(self: any) -> String
}

impl Sendable for EmailAlert {
    fn send(self: any) -> String {
        return "Sending email to " + self.to + ": " + self.subject
    }
}

impl Loggable for EmailAlert {
    fn log_entry(self: any) -> String {
        return "[EMAIL] to=" + self.to + " subject=" + self.subject
    }
}

let alert = EmailAlert {
    to: "ops@example.com",
    subject: "Server down",
    body: "The production server is unreachable."
}

print(alert.send())      // Sending email to ops@example.com: Server down
print(alert.log_entry()) // [EMAIL] to=ops@example.com subject=Server down
\end{lstlisting}

Each trait captures one aspect of behaviour. A \lstinline|SlackAlert| struct
could implement the same traits with different bodies, and any code that
calls \lstinline|.send()| or \lstinline|.log_entry()| would work with
either type.

\subsection{Naming Conventions}

Good trait names describe a capability, not a thing:

\begin{itemize}
  \item \lstinline|Describable|, \lstinline|Displayable| --- ``can be described / displayed''
  \item \lstinline|HasArea|, \lstinline|HasLength| --- ``possesses a measurable property''
  \item \lstinline|Scalable|, \lstinline|Combinable| --- ``can be scaled / combined''
  \item \lstinline|Collection|, \lstinline|Geometry| --- group of related operations
\end{itemize}

When you read \lstinline|impl Scalable for Vec2|, the intent is
immediately clear: ``\lstinline|Vec2| can be scaled.''

%% ─────────────────────────────────────────────────────────────────────
\section{Trait-Driven Polymorphism}\label{sec:trait-polymorphism}
\index{polymorphism}
\index{trait!polymorphism}

Because Lattice is dynamically typed, any value can be passed to any
function. Traits add a \emph{semantic} layer on top: they document which
methods a value must support, and the runtime dispatches to the correct
implementation based on the value's type.

\subsection{Heterogeneous Collections}

You can put different types in the same array and call the same method on
each:

\begin{lstlisting}[language=Lattice, caption={Polymorphic dispatch over an array}]
struct Circle { radius: Float }
struct Rect { width: Float, height: Float }

trait HasArea {
    fn area(self: any) -> Float
}

impl HasArea for Circle {
    fn area(self: any) -> Float {
        return 3.14159 * self.radius * self.radius
    }
}

impl HasArea for Rect {
    fn area(self: any) -> Float {
        return self.width * self.height
    }
}

let shapes = [
    Circle { radius: 3.0 },
    Rect { width: 4.0, height: 5.0 },
    Circle { radius: 1.0 }
]

flux total_area = 0.0
for shape in shapes {
    total_area = total_area + shape.area()
}
print(total_area)  // 28.27431 + 20.0 + 3.14159 = 51.4159
\end{lstlisting}

The \lstinline|for| loop does not know whether each element is a
\lstinline|Circle| or a \lstinline|Rect|. It calls \lstinline|.area()|
on each, and the runtime dispatches to the correct implementation based
on the struct's type name.

\begin{defnbox}{Dynamic Dispatch}
Lattice uses \emph{dynamic dispatch} for trait methods. When you call
\lstinline|shape.area()|, the VM looks up the method key
\texttt{TypeName::area} at runtime, using the actual type of the value.
This is different from static dispatch (resolved at compile time) used by
languages like Rust.
\end{defnbox}

\subsection{Methods in Expressions and Conditionals}

Trait method calls are expressions---you can use them in \lstinline|if|
conditions, arithmetic, and anywhere else a value is expected:

\begin{lstlisting}[language=Lattice, caption={Methods in expressions}]
struct Rect { w: any, h: any }

trait Geometry {
    fn area(self: any) -> any
}

impl Geometry for Rect {
    fn area(self: any) -> any {
        return self.w * self.h
    }
}

let r1 = Rect { w: 4, h: 6 }
let r2 = Rect { w: 3, h: 8 }

let total_area = r1.area() + r2.area()
print(total_area)  // 48

let larger = if r1.area() > r2.area() { "r1" } else { "r2" }
print(larger)  // r2
\end{lstlisting}

\subsection{Methods in Loops}

A common pattern is to accumulate results by calling trait methods inside
a loop:

\begin{lstlisting}[language=Lattice, caption={Accumulating over trait methods}]
struct Rect { w: any, h: any }

trait Geometry {
    fn area(self: any) -> any
}

impl Geometry for Rect {
    fn area(self: any) -> any {
        return self.w * self.h
    }
}

let rects = [
    Rect { w: 1, h: 1 },
    Rect { w: 2, h: 2 },
    Rect { w: 3, h: 3 }
]

flux total = 0
for r in rects {
    total = total + r.area()
}
print(total)  // 1 + 4 + 9 = 14
\end{lstlisting}

\subsection{Traits on Enums}

Traits are not limited to structs---you can implement them for enums too:

\begin{lstlisting}[language=Lattice, caption={A trait implemented for an enum}]
enum Status {
    Active,
    Inactive,
    Error(any)
}

trait Displayable {
    fn display(self: any) -> String
}

impl Displayable for Status {
    fn display(self: any) -> String {
        return "Status:active"
    }
}

let s = Status::Active
print(s.display())  // Status:active
\end{lstlisting}

\begin{notebox}{Enum Method Dispatch}
When you call a trait method on an enum value, the runtime looks up the
method key using the \emph{enum type name}---not the variant name. So
\lstinline|Status::Active|, \lstinline|Status::Inactive|, and
\lstinline|Status::Error("oops")| all dispatch to the same impl block.
If you need per-variant behaviour, use \lstinline|match| inside the
method body.
\end{notebox}

\subsection{A Complete Polymorphism Example}

Let's put it all together with a two-dimensional vector type that
supports scaling and combining:

\begin{lstlisting}[language=Lattice, caption={A complete trait-driven design}]
struct Vec2 { x: any, y: any }

trait Describable {
    fn describe(self: any) -> String
}

trait Scalable {
    fn scale(self: any, factor: any) -> any
}

trait Combinable {
    fn combine(self: any, other: any) -> any
}

impl Describable for Vec2 {
    fn describe(self: any) -> String {
        return "Vec2(" + to_string(self.x) + ", " + to_string(self.y) + ")"
    }
}

impl Scalable for Vec2 {
    fn scale(self: any, factor: any) -> any {
        return Vec2 { x: self.x * factor, y: self.y * factor }
    }
}

impl Combinable for Vec2 {
    fn combine(self: any, other: any) -> any {
        return Vec2 { x: self.x + other.x, y: self.y + other.y }
    }
}

let velocity = Vec2 { x: 3, y: 4 }
let boost = velocity.scale(2)
print(boost.describe())  // Vec2(6, 8)

let gravity = Vec2 { x: 0, y: -1 }
let result = boost.combine(gravity)
print(result.describe())  // Vec2(6, 7)
\end{lstlisting}

This design is extensible. If you later add a \lstinline|Vec3| struct, you
can implement the same three traits for it, and any code that operates on
\lstinline|Describable|, \lstinline|Scalable|, or \lstinline|Combinable|
values will work with both types---no changes needed.

%% ─────────────────────────────────────────────────────────────────────
\section{Under the Hood}\label{sec:trait-internals}
\index{trait!internals}

Understanding the implementation helps when debugging method dispatch
or designing complex trait hierarchies.

\subsection{Compilation}

Traits themselves are metadata-only. The compiler sees a trait declaration
and records it, but emits no bytecode for it. Only \lstinline|impl| blocks
produce code:

\begin{enumerate}
  \item Each method body is compiled as a regular function.
  \item The compiled function is registered as a global under the key
    \texttt{"TypeName::methodName"}.
  \item The VM resolves this key at call time via the \texttt{OP\_INVOKE}
    opcode family.
\end{enumerate}

You can see this in \filename{src/stackcompiler.c}---the
\texttt{ITEM\_TRAIT} case is essentially a no-op, while \texttt{ITEM\_IMPL}
compiles each method and emits \texttt{OP\_DEFINE\_GLOBAL}.

\subsection{Dispatch Performance}

The stack VM uses three specialised invoke opcodes for performance:

\begin{itemize}
  \item \textbf{\texttt{OP\_INVOKE\_LOCAL}}: When the receiver is a local
    variable, the VM can skip the stack and access the value directly from
    the frame's slot array. This is the fastest path.
  \item \textbf{\texttt{OP\_INVOKE\_GLOBAL}}: When the receiver is a global
    variable. Slightly slower due to the global lookup.
  \item \textbf{\texttt{OP\_INVOKE}}: The general case for arbitrary
    expressions. The struct is on the stack, and the VM searches for the
    method.
\end{itemize}

Each path also maintains a \emph{polymorphic inline cache} (PIC) to avoid
repeated lookups when the same call site sees the same type across
iterations. You can explore this in \filename{src/stackvm.c}.

\subsection{Self Cloning}

When a method is invoked, the runtime clones \lstinline|self| via
\texttt{value\_deep\_clone()} before passing it to the method frame.
This ensures that mutations inside the method cannot accidentally affect
the caller's copy---maintaining Lattice's value semantics guarantee.

%% ─────────────────────────────────────────────────────────────────────
\section{Exercises}\label{sec:trait-exercises}

\begin{enumerate}
  \item \textbf{Shape hierarchy.}
    Define structs \lstinline|Circle|, \lstinline|Square|, and
    \lstinline|Triangle|. Create a \lstinline|HasArea| trait and implement
    it for all three. Put instances of each into an array and compute the
    total area in a loop.

  \item \textbf{Stack from scratch.}
    Using the \lstinline|Collection| trait from this chapter (with methods
    \lstinline|push_item|, \lstinline|pop_item|, \lstinline|peek_item|,
    and \lstinline|is_empty|), implement a \lstinline|Queue| struct that
    supports FIFO ordering. Verify that items come out in insertion order.

  \item \textbf{Composing traits.}
    Define two traits: \lstinline|Serializable| (with a
    \lstinline|to_string_repr| method) and \lstinline|Parseable| (with a
    class-level \lstinline|from_string| function). Implement both for a
    \lstinline|Temperature| struct that holds a value and a unit
    (\lstinline|"C"| or \lstinline|"F"|). Demonstrate round-tripping:
    convert to a string and parse back.

  \item \textbf{Trait method chaining.}
    Create a \lstinline|Builder| struct with fields \lstinline|parts| and
    \lstinline|count|. Implement a \lstinline|Buildable| trait with methods
    \lstinline|add_part| (returns a new Builder with the part appended) and
    \lstinline|build| (returns a summary string). Chain at least four
    \lstinline|add_part| calls followed by \lstinline|build|.
\end{enumerate}

%% ─────────────────────────────────────────────────────────────────────
\section*{What's Next}

Traits let you write code that works across multiple types. But so far,
every type annotation in our examples has been \lstinline|any|---a catch-all
that says nothing about what the value actually is. In the next chapter we
meet \emph{generics}, which let you parameterise functions, structs, and
enums over types. Instead of \lstinline|fn identity(x: any) -> any|, you
will write \lstinline|fn identity<T>(x: T) -> T|---and the language will
understand the relationship between input and output.
